\documentclass[11pt]{article}

\usepackage[margin=1in]{geometry}
\usepackage[T1]{fontenc}
\usepackage[utf8]{inputenc}
\usepackage{lmodern}
\usepackage{amsmath,amssymb}
\usepackage{booktabs}
\usepackage{array}
\usepackage{enumitem}
\usepackage{longtable}
\usepackage{xcolor}
\usepackage{hyperref}
\usepackage{listings}

\hypersetup{
    colorlinks=true,
    linkcolor=blue,
    urlcolor=blue
}

\lstset{
    basicstyle=\ttfamily\small,
    breaklines=true,
    columns=fullflexible,
    keepspaces=true,
    frame=single
}

\title{Campus Image Branches}
\author{}
\date{}

\begin{document}
\maketitle

\section{Semantic Color Contract}

The campus maps are not ordinary thresholded images. Their pixel colors encode semantics that affect:
\begin{itemize}[leftmargin=2em]
    \item traversability,
    \item spawn eligibility,
    \item zone membership,
    \item dynamic-obstacle eligibility for the dynamic campus branch.
\end{itemize}

The current campus loader recognizes:
\begin{itemize}[leftmargin=2em]
    \item colored zones as traversable and spawnable,
    \item white walkways as traversable but non-spawnable,
    \item gray regions as non-traversable.
\end{itemize}

\section{Current Campus Branch Split}

The project now uses the two campus source images in the following way:
\begin{itemize}[leftmargin=2em]
    \item \texttt{static\_campus\_area\_1}: static image branch,
    \item \texttt{dynamic\_campus\_area\_2}: dynamic image branch.
\end{itemize}

Each campus branch now also has a \texttt{narrow\_lanes} boolean in \texttt{master\_config.py}. That toggle selects between:
\begin{itemize}[leftmargin=2em]
    \item \texttt{dynamic\_campus\_area\_1/campus\_area\_1\_narrow\_lanes.png} or \texttt{campus\_area\_1\_wide\_lanes.png},
    \item \texttt{dynamic\_campus\_area\_2/campus\_area\_2\_narrow\_lanes.png} or \texttt{campus\_area\_2\_wide\_lanes.png}.
\end{itemize}

The branch roles themselves do not change when the lane variant changes: \texttt{static\_campus\_area\_1} remains the static campus branch and \texttt{dynamic\_campus\_area\_2} remains the dynamic campus branch.

\section{\texttt{static\_campus\_area\_1}}

Campus area 1 is now the static campus branch. In branch-role terms, this is the branch that uses the two-zone campus layout. It samples:
\begin{itemize}[leftmargin=2em]
    \item a dispersed start group from one zone,
    \item a clustered goal group from the other zone,
    \item then assigns those starts and goals one-to-one.
\end{itemize}

The start zone and goal zone can swap across runs, but they must remain different.

\section{\texttt{dynamic\_campus\_area\_2}}

Campus area 2 is now the dynamic campus branch. In branch-role terms, this is the branch that uses the three-zone campus layout. It samples:
\begin{itemize}[leftmargin=2em]
    \item a clustered start group from one zone,
    \item a clustered goal group from exactly one \emph{different} zone,
    \item then assigns those starts and goals one-to-one.
\end{itemize}

So if starts are in zone A, goals must be entirely in zone B or entirely in zone C, not split across both.

\section{Dynamic Generation Policy}

For \texttt{dynamic\_campus\_area\_2}, dynamic obstacle generation still uses:
\[
\texttt{dynamic\_generation\_cell\_mode=zone\_colors\_only}.
\]
That means dynamic obstacles are generated only inside the designated zone colors. The static image layout itself is preserved.

\section{Sampling Policy}

The current campus assignment logic now follows these rules:
\begin{itemize}[leftmargin=2em]
    \item one goal per agent,
    \item starts and goals are unique,
    \item dispersed sets keep internal 8-neighbor separation while clustered sets follow the global \texttt{compact\_clustering} switch,
    \item starts and goals may be adjacent to one another,
    \item when campus distinct-zone mode is active, the start zone and goal zone must differ,
    \item optional individual reachability is checked before accepting a sample.
\end{itemize}

\section{Renderer Note}

The campus preparation code stores gray-region vertices as \texttt{visually\_free\_vertices}. During visualization regeneration, the renderer now also builds a render-only composite-map copy that temporarily converts those gray obstacle vertices into free-space cells before drawing. In practice, this makes gray campus regions appear as ordinary white open areas in Pillow outputs even though they remain non-traversable for the actual solver and are still excluded from sampling.

This behavior is now enforced both for freshly computed runs and for visualization-only reruns that reuse already persisted raw MAPF data. During static-campus visualization regeneration, the renderer reconstructs the gray-cell mask from the \emph{current} campus image and applies it as a render-only override, so old persisted run data can still be visualized correctly without recomputing the solver stage.

\section{Practical Warning}

If a campus run logs something like \texttt{zone\_cells=0} or otherwise behaves as though no campus zones exist, the first thing to check is the source image itself. The campus branch assumes the image was actually recolored with the intended zone colors. If the image colors are wrong, the semantic loader cannot recover the intended zones. Also verify that \texttt{master\_config.py} still points to the intended branch/image pairing and the intended \texttt{narrow\_lanes} variant.

\end{document}
